\documentclass[12pt,a4paper]{report}
\usepackage[utf8]{inputenc}
\usepackage{amsmath}
\usepackage{amsfonts}
\usepackage{amssymb}
\usepackage{amsthm}
\usepackage{hyperref}

\usepackage{multicol}
\usepackage{fancyhdr}
\usepackage[inline]{enumitem}
\usepackage{tikz}
\usepackage{tikz-cd}
\usetikzlibrary{calc}
\usetikzlibrary{shapes.geometric}
\usetikzlibrary{positioning}
\usepackage[margin=0.5in]{geometry}
\usepackage{xcolor}

\hypersetup{
    colorlinks=true,
    linkcolor=blue,
    filecolor=magenta,      
    urlcolor=cyan,
    pdftitle={Tensors},
    pdfpagemode=FullScreen,
    }

%\urlstyle{same}

\newcommand{\CLASSNAME}{Math 5050 -- Special Topics: Tensor Analysis}
\newcommand{\STUDENTNAME}{Paul Carmody}
\newcommand{\ASSIGNMENT}{Chapter 13: Embedded Curves}
\newcommand{\DUEDATE}{March 2026}
\newcommand{\PROFESSOR}{Professor Berchenko-Kogan}
\newcommand{\SEMESTER}{Spring 2026}
\newcommand{\SCHEDULE}{TBD}
\newcommand{\ROOM}{Remote}

\newcommand{\MMN}{M_{m\times n}}
\newcommand{\FF}{\mathcal{F}}

\pagestyle{fancy}
\fancyhf{}
\chead{ \fancyplain{}{\CLASSNAME} }
%\chead{ \fancyplain{}{\STUDENTNAME} }
\rhead{\thepage}

\fancyfoot{} % clear all footer fields
\fancyfoot[LE,RO]{\thepage}
\fancyfoot[LO,CE]{-------\\\ASSIGNMENT}
%\fancyfoot[CO,RE]{To: Dean A. Smith}
\input{../MyMacros}

\newcommand{\RED}[1]{\textcolor{red}{#1}}
\newcommand{\BLUE}[1]{\textcolor{blue}{#1}}

\newcommand{\SUPP}{\operatorname{supp}}
\newcommand{\CLOSURE}{\operatorname{closure of}}
\newcommand{\BD}{\operatorname{bd}}
\newcommand{\HHN}{\mathcal{H}^n}
%\newcommand{\COFF}[3]{\Gamma^{#1}_{{#2}{#3}}}
\newcommand{\COFF}[3]{\PART{\textbf{Z}_{#2}}{Z^{#3}}\cdot\textbf{Z}^{#1}}
\newcommand{\COFFF}[3]{\frac{1}{2}U^{{#1}\Omega}\PAREN{\PART{U_{\Omega{#2}}}{U^{#3}} + \PART{U_{\Omega{#3}}}{U^{#2}} - \PART{U_{{#2}{#3}}}{U^\Omega}}}
\newcommand{\VK}[1]{\textbf{#1}}
\newcommand{\VKR}{\VK{R}}
\newcommand{\VKZ}{\VK{Z}}

\newcommand{\RIEMANN}[4]{ \PART{\Gamma^#1_{#4#2}}{S^#3} - \PART{\Gamma^#1_{#3#2}}{S^#4} + \Gamma^#1_{#3\omega}\Gamma^\omega_{#4#2} - \Gamma^#1_{#4\omega}\Gamma^{\omega}_{#3#2}}

\begin{document}

\begin{center}
	\Large{\CLASSNAME -- \SEMESTER} \\
	\large{ w/\PROFESSOR}
\end{center}
\begin{center}
	\STUDENTNAME \\
	\ASSIGNMENT -- \DUEDATE\\
\end{center} 

\textbf{\Large{Notes:}}

\begin{description}
	\item \textbf{Elements that make up a tensor-space.}

	The method is \textbf{First} to define the space and transitions between it and the ambient space (i.e., space aligned with the scalar field).  \textbf{Next}, define shift tensor $\to$ covariant tensor $\to$ covariant metric tensor $\to$ contravariant metric tensor $\to$ contravariant tensor.  This is the landscape of the space.  \textbf{Next}, define the Levi-Civita and Christoffel symbols $\to$ covarariant derivative $\to$ contravariant derivative.  This is the landscape of change over the space.  \textbf{Finally}, define the curvature tensor $\to$ curvature mean.

	\begin{itemize}
		\item \textbf{Define the space.}  

All of this academic but establishes the foundation for embedded surfaces, hyper-surfaces and so forth.  The space is represented by $U$ and its dimension can only be 1.  However, to make the proper analogies to multi-dimensional embedded spaces, we will refer to the coordinates as $\Phi, \Theta,$ and $\Psi$ all of which will be 1.

		\item \textbf{Establish link to ambient space. Covariant vectors, and shift tensor.}

		Thus, $Z^i = Z^i(U)$.  Here, $U$ is a point in our space and $Z^i$ represent the coordinates generated by the scalar field.  That is, basis generated by the gradient.  Therefore the shift tensor will be directly
			\begin{align*}
				Z^i_\Phi &= \PART{Z^i}{U_\Phi} & \text{shift tensor} \\
				\textbf{U}_\Phi &= \PART{\textbf{R}(U)}{U^\Phi} = \textbf{Z}_iZ^i_\Phi & \text{covariant tensor} \\
			\end{align*}\BLUE{Interpretaion: $\textbf{U}_\Phi$ can only be one vector, thus the expression $\textbf{Z}_iZ^i_\Phi $ maps all $ \textbf{Z}_i$ to $\textbf{U}_\Phi$ }

		\item \textbf{Establish a covariant metric tensor and contravaiant metric tensor which leads to the contravaiant vector.  Along with a `volume' tensor.}

			\begin{align}
				U_{\Phi\Psi} &= \textbf{U}_\Phi\cdot \textbf{U}_\Psi = Z_{ij}Z^i_\Phi Z^j_\Psi & \text{covariant metric tensor}\\
				U_{\Theta\Phi}U^{\Phi\Psi} &= \delta^\Psi_\Theta  & \text{contravariant metric tensor definition}\\
				\textbf{U}^\Phi &= U^{\Phi\Psi}\textbf{U}_\Psi & \text{covariant vector}\\
				U &= |U_{..}| & \text{volume/length element} \\
				&= Z_{ij}\PART{Z^i}{U^1}\PART{Z^j}{U^1}
			\end{align}\BLUE{Interpretation: (1) ???? (2) the sum of one term where $\Phi,\Psi$ can be only 1, (3) all this does is reinterpet $\textbf{U}$ as a contravariant vector, (4) reinterpret $U$ as a determinant of the metric, (5) and see its connection to the ambient space.}
			
			We note that 
			\begin{align*}
				U_{\Phi\Psi} &= U^{\Phi\Psi} = [1] & \text{ a $1\times 1$ matrix}\\
				\Gamma^\Theta_{\Phi\Psi} &=0 & \text{defined below}
			\end{align*}

		\item \textbf{Establish the Levi-Civita symbol and Christoffel symbol.}

		By analogy the Levi-Civita symbol becomes
		\begin{align*}
			\varepsilon_\Phi &= \sqrt{U} \\
			\varepsilon^\Phi &= \frac{1}{\sqrt{U}} 
		\end{align*} and by definition the Christoffel Symbol becomes
		\begin{align*}
			\Gamma^\Theta_{\Phi\Psi} &= \textbf{U}^\Theta \cdot \PART{\textbf{U}_\Phi}{U^\Psi} = -\PART{\textbf{U}^\Phi}{U^\Psi} \cdot \textbf{U}_\Theta  \\
				&= \COFF{\Theta}{\Phi}{\Psi}
		\end{align*}

		\item \textbf{Establish the covariant derivative.}

			For any tensor $W^{i..}_{j..}$ we can define the covariant derivative using the shift tensor and ambient covariant derivative as
			\begin{align*}
				\nabla_\Theta W^{i..}_{j..} &= Z^k_\theta \nabla_k W^{i..}_{j..}
			\end{align*}Since we only have one dimension and it refers to the arc length we also have
			\begin{align*}
				\nabla_s = \varepsilon^\Theta \nabla_\Theta
			\end{align*}

			\BLUE{Interpretation: up until this point, we've equated all of the necessary ingredients to understand tensors on a manifold/space.  In this case, a curve on a surface.  Each 'term' (in a literary sense) has more or less the same value but not the same meaning.  Furthermore, there is a tensor algebraic necessity to convert '1' in one form into '1' into another (i.e, a vector of length 1 to a metric of size 1).  The lesson here is to establish and understand how the tensor framework defines a scalar field in greater detail.}

		\item \textbf{Establish the curvature tensor.}

			\begin{align*}
				\textbf{B}_{\Phi\Psi} &= \nabla_\Phi\textbf{U}_\Psi & \text{\textbf{Curvature tensor, a single vector}.} \\
				\textbf{B}^\Phi_\Phi &= \nabla_s^2\textbf{R} = \frac{d^2\textbf{R}(s)}{dx^2} & \text{curvature normal, $1\times 1$ matrix} \\
				 &= \kappa \textbf{P} & \text{where $\textbf{P}$ is called the \DEFINE{Principle Normal}}
			\end{align*}

		\item \textbf{Normals and tangents to a curve in space.}

		The Unit Tangent vector and Principle Normal are defined/related by
			\begin{align*}
				\frac{d\textbf{R}}{ds} &= \textbf{T} \\
				\frac{d\textbf{T}}{ds} &= \kappa \textbf{P} \\
				\textbf{P}\cdot\textbf{P} &= 1 \\
				\textbf{P}\cdot\textbf{T} &= 0
			\end{align*}
	\end{itemize}
\end{description}

\HLINE

\begin{description}

\section*{Section 13.2}
\item \textbf{Exercise 277} Explain why $\varepsilon^\Phi\textbf{U}_\Phi$ is the unit tangent.  Hint: Consider $U^1=s$.

\section*{Section 13.5}
\item \textbf{Exercise 278} Show that $\nabla_s$ coincides with $d/ds$ for invariants
\begin{align*}
	\nabla_s W = \frac{dW}{ds}
\end{align*}

\item \textbf{Exercise 279}  Show that $\nabla_s$ satisfies the sum and product rules.

\item \textbf{Exercise 280}  Show that $\nabla_s$ commutes with contraction.

\BLUE{WTS: $\nabla_s ($contraction$)=$ contraction$(\nabla_s)$.  That is, given a tensor $T^i_j$ and $S^i_{sj} = \nabla_s T^i_j$ then $\nabla_s T^i_i = S^i_{ij}$.
}

\item \textbf{Exercise 281}  Show that $\nabla_s$ satisfies the following chain rule for curve retrictions of ambient variants
\begin{align*}
	\nabla_s W^i_j=\varepsilon^\Theta Z^k_\Theta \nabla_k W^i_j
\end{align*}Therefore, 
\begin{align*}
	\nabla_s W^i_j = T^k\nabla_k W^i_j
\end{align*}where $T_k$ is the contravariant component of the unit tangent.

\item \textbf{Exercise 282} Show that $$\nabla_s^2= \nabla_\Phi \nabla^\Phi$$.

\section*{Section 13.7}

\item \textbf{Exercise 283} Differentiate the identity $\textbf{T}(s)\cdot\textbf{T}(s)=1$ to conclude that $\textbf{P}$ is orthogonal to $\textbf{T}$.

\BLUE{
\begin{align*}
	\frac{d}{ds}(\textbf{T}(s)\cdot\textbf{T}(s)) &= \frac{d}{ds}(\textbf{T}(s))\cdot\textbf{T}(s) + \textbf{T}(s)\cdot\frac{d}{ds}(\textbf{T}(s)) \\
		&= \kappa\textbf{P}\cdot\textbf{T} + \textbf{T}\cdot(\kappa\textbf{P}) \\
		&= \kappa\textbf{P}\cdot\textbf{T} \\
	\therefore \textbf{P}\cdot\textbf{T} &= 0 \implies \textbf{P} \perp \textbf{T}
\end{align*}
}

\item \textbf{Exercise 284} Show that the vector $d\textbf{P}/ds + \kappa\textbf{T}$ is obtained from $d\textbf{P}/ds$ by the Gram-Schmidt orthogonalization procedure.

\item \textbf{Exercise 285} Calculate the principal normal, binormal, curvature, torsion for the spiral
\begin{align*}
	x(t) &= 4\cos t \\
	y(t) &= 4\sin t \\
	z(t) &= 3t.
\end{align*}

\newcommand{\DT}[1]{\frac{d{#1}}{dt}}
\BLUE{
\begin{itemize}
	\item Determine $\textbf{T}$.
		\begin{align*}
			\textbf{T} &= \frac{\textbf{R}}{dt} \\
				&= \PAREN{\DT{x}, \DT{y}, \DT{z}} \\
				&= \PAREN{ -4\sin t, 4\cos t, 3}
		\end{align*}
	\item Determine $\textbf{P}$.
		\begin{align*}
			\kappa\textbf{P} &= \DT{\textbf{T}} \\
				&= \PAREN{-4\cos t, -4\sin t, 0} \\
			\textbf{P} &= \PAREN{\cos t, \sin t, 0}
		\end{align*}
	\item Determine $\textbf{Q}$.
		\begin{align*}
			\tau\textbf{Q} &= \DT{\textbf{P}}+\kappa\textbf{T} \\
				&= ( -\sin t, \cos t, 0) + -4 \PAREN{ -4\sin t, 4\cos t, 3} \\
				&= (15\sin t, -15\cos t, -12)
		\end{align*}
	\item Determine $\kappa$.\\
		From $\textbf{P}$ we can see that $\kappa = -4$.
	\item Determine $\tau$.
		From $\textbf{Q}$ we can see that $\tau = 15$.
\end{itemize}
}

\item \textbf{Exercise 286} Repeat the calculation for the spiral
\begin{align*}
	x(t) &= 4\cos t \\
	y(t) &= 4\sin t \\
	z(t) &= -3t.
\end{align*}Not which answers change for this spiral.

\section*{Section 13.8}

\item \textbf{Exercise 287}  Explain why the last equation in the Frenet procedure reads
\begin{align*}
	\frac{d\textbf{T}_{N-1}}{ds} &= \kappa_{N-1}\textbf{T}_{N-2}
\end{align*}consistent with equation (13.57).

\begin{align*}
	\frac{d\textbf{Q}}{ds} &= -\tau\textbf{P} & (13.57)
\end{align*}

\BLUE{Let $T_0 = T$ from the \textit{Frenet Formulas} (13.55, 13.56, 13.57).  Then, $P = T_1$ and $Q = T_2$.  And, from (13.69)
\begin{align*}
	\frac{\textbf{T}_m}{ds}\cdot\textbf{T}_{m-1} &= -\kappa_m & (13.69)
\end{align*}we can define $\kappa = \kappa_1$ and $\tau = \kappa_2$. Also, all $\kappa_i$ are the dot products of other unit vectors, which have all now be accounted for.  Thus, the next level or processing would make $\textbf{T}_4$ perpendicular to the three other vectors, thus $\kappa_3 = 0$.
\begin{align*}
	\frac{d\textbf{T}_2}{ds} &= \kappa_{1}\textbf{T}_0
\end{align*}
}

\section*{Section 14.7}
\item \textbf{Exercise 288}  Define the unit tangent vector $\textbf{T}$
\begin{align*}
	T^i = \varepsilon_{ijk}N^jn^\alpha z^k_\alpha.
\end{align*}
Show that the vector $\textbf{T}$ with components give by (14.43) is inded the unit tangent to the contour boundary.  That is, $\textbf{T}$ is unit length
\begin{align*}
	\textbf{T}\cdot\textbf{T} = 1
\end{align*}orthogonal to the surface normal $\textbf{N}$
\begin{align*}
	\textbf{T}\cdot\textbf{N} = 0
\end{align*}and the countour normal
\begin{align*}
	\textbf{T}\cdot\textbf{n} = 0.
\end{align*}

\BLUE{Converting to ambiant coordinates we get
\begin{align*}
	n^i &= n^\alpha Z^i_\alpha \\
	T^i &= \varepsilon_{ijk}N^jn^k \\
	\textbf{T} &= \textbf{N}\times\textbf{n}
\end{align*}Take the inner product
\begin{align*}
	\textbf{T}\cdot\textbf{N} &= (\textbf{N}\times\textbf{n})\cdot\textbf{N} = 0\\
	\textbf{T}\cdot\textbf{n} &= (\textbf{N}\times\textbf{n})\cdot\textbf{n} = 0
\end{align*}Since $|\textbf{N}| = |\textbf{n}| = 1$ then $|\textbf{T}|=1$.
}

\item \textbf{Exercise 289} Explain why the term $\varepsilon_{ijk}Z^j_\beta Z^k_\alpha B^{\alpha\beta}$.

\end{description}
\end{document}
